% capitolo 5 - Installazione e configurazione + data ingestion dei primi dataset
\chapter{Installazione e configurazione}
\label{cap:installazione}

\intro{In questo capitolo vengono forniti i prerequisiti e le indicazioni per l'installazione e il \emph{setup} di OpenMetadata in ambiente locale.}\\

\section{Prerequisiti}
L'installazione e l'esecuzione di OpenMetadata prevedono requisiti specifici, legati alle risorse hardware della macchina ospitante e 
all'ambiente di esecuzione locale, basato in questo contesto sulla containerizzazione tramite Docker.

Per garantire il corretto funzionamento di una distribuzione locale di test, la macchina ospitante (basata su sistema operativo Linux oppure su Windows 
tramite \gls{wsl}2 o macchina virtuale) deve allocare all'\emph{engine} di Docker le seguenti risorse hardware minime:
\begin{itemize}
    \item \textbf{CPU}: Almeno 4 vCPU.
    \item \textbf{Memoria RAM}: Almeno 6--8 GB di RAM interamente dedicati ai container.
    \item \textbf{Spazio su disco}: Almeno 30--40 GB di spazio libero.
\end{itemize}

Dal punto di vista software, i requisiti fondamentali prevedono Git per il controllo di versione, Docker Engine (versione 20.10.0 o superiore) e Docker Compose (versione v2.1.1 o superiore). 
L'architettura locale di OpenMetadata viene orchestrata attraverso quattro container distinti:
\begin{enumerate}
    \item \textbf{Database di sistema}: Un'istanza PostgreSQL o MySQL dedicata al salvataggio dello stato, delle configurazioni e dei metadati gestiti dalla piattaforma.
    \item \textbf{OpenMetadata Server}: Il backend centrale della piattaforma sviluppato in Java.
    \item \textbf{OpenMetadata Ingestion (\emph{Apache Airflow})}: Il motore di orchestrazione responsabile dell'esecuzione automatica e pianificata delle pipeline di estrazione dati.
    \item \textbf{Motore di ricerca}: Un container ElasticSearch (o \gls{opensearch}) dedicato all'indicizzazione dei metadati e alla ricerca \emph{full-text}.
\end{enumerate}

Per consentire la corretta fruizione dei servizi sul browser, l'ambiente di rete locale deve garantire che le porte 8585 (dedicata all'interfaccia utente di OpenMetadata) 
e 8080 (dedicata al pannello di controllo di Apache Airflow) siano libere da conflitti.

Infine, per testare la piattaforma senza violare i vincoli di riservatezza sui dati reali aziendali, durante lo stage è stato utilizzato il 
database relazionale di esempio \emph{AdventureWorks} fornito da Microsoft\cite{site:adventureworks}. 
Tale risorsa è stata ripristinata all'interno di un'istanza locale di Microsoft SQL Server, simulando uno scenario di produzione realistico, completo di 
relazioni e chiavi esterne su cui eseguire le prime procedure di \emph{Data Ingestion} e \emph{Data Discovery}.

\section{Installazione di OpenMetadata}
Una volta soddisfatti tutti i prerequisiti, la procedura di deployment locale di OpenMetadata è stata eseguita tramite Docker e Docker Compose, 
seguendo i passaggi della documentazione ufficiale\cite{site:openmetadata-docs}, di seguito elencati:
\begin{enumerate}
    \item \textbf{Creazione della directory di lavoro}: è stata creata una cartella dedicata al progetto:
    \begin{verbatim}
mkdir openmetadata-docker && cd openmetadata-docker
    \end{verbatim}
    \item \textbf{Download del file di configurazione}: è stato scaricato il file \texttt{docker-compose.yml} ufficiale dal repository\cite{site:openmetadata-repo} del progetto (oppure il file \texttt{docker-compose-postgres.yml} qualora si preferisca utilizzare PostgreSQL come database di sistema al posto di MySQL).
    \item \textbf{Avvio dei servizi}: i container Docker di OpenMetadata sono stati eseguiti tramite il seguente comando:
    \begin{verbatim}
docker compose -f docker-compose.yml up --detach
    \end{verbatim}
\end{enumerate}

Una volta che i container risultano operativi, la piattaforma è pronta per l'utilizzo:
\begin{itemize}
    \item \textbf{Interfaccia utente di OpenMetadata}: raggiungibile da browser all'indirizzo \nolinkurl{http://localhost:8585}. Al primo accesso, l'autenticazione avviene tramite le credenziali di default fornite dal sistema: \texttt{admin@open-metadata.org} come \emph{username} e \texttt{admin} come \emph{password}. Da qui è poi possibile accedere alle impostazioni generali e configurare nuovi utenti;
    \item \textbf{Pannello di controllo di Apache Airflow}: raggiungibile all'indirizzo \nolinkurl{http://localhost:8080} inserendo le credenziali di default (\texttt{admin} / \texttt{admin}). Airflow opera in background come motore di orchestrazione per le pipeline di \emph{ingestion} dei metadati; sebbene la gestione delle sincronizzazioni avvenga interamente dall'interfaccia di OpenMetadata, l'accesso diretto ad Airflow risulta fondamentale per l'analisi dei log e il \emph{debugging} avanzato in caso di fallimento delle pipeline.
\end{itemize}


\section{Data Ingestion}

\section{Ricerca e Data Discovery}

\section{Data Lineage}